% Термины и определения
\newglossaryentry{id4}{
    name={Адрес},
    description={Уникальный идентификатор в сети Bitcoin, используемый для отправки и получения средств}
}

\newglossaryentry{id15}{
    name={Батч},
    description={Подмножество обучающих данных, используемое для одного шага обновления весов нейронной сети}
}

\newglossaryentry{id2}{
    name={Блокчейн},
    description={Распределённая база данных, состоящая из последовательно связанных блоков, содержащих информацию о транзакциях}
}

\newglossaryentry{id10}{
    name={Гиперпараметр},
    description={Параметр модели или процесса обучения, задаваемый до начала обучения и не изменяемый в его ходе}
}

\newglossaryentry{id20}{
    name={Граф транзакций},
    description={Графовое представление Bitcoin-сети, где узлами являются адреса, а рёбрами~--- транзакции между ними}
}

\newglossaryentry{id7}{
    name={Датасет},
    description={Набор данных, используемый для обучения, тестирования и валидации алгоритмов машинного обучения}
}

\newglossaryentry{id14}{
    name={Дропаут},
    description={Метод регуляризации нейронных сетей, при котором случайные нейроны отключаются во время обучения для предотвращения переобучения}
}

\newglossaryentry{id19}{
    name={Извлечение признаков},
    description={Процесс выделения характеристик из данных для использования в алгоритмах машинного обучения}
}

\newglossaryentry{id13}{
    name={Классификация},
    description={Задача машинного обучения, заключающаяся в отнесении объектов к предопределённым классам}
}

\newglossaryentry{id11}{
    name={Кросс-валидация},
    description={Метод оценки качества модели путём разбиения данных на несколько частей и поочерёдного использования каждой из них в качестве тестовой}
}

\newglossaryentry{id12}{
    name={Матрица ошибок},
    description={Таблица, отображающая соотношение предсказанных и истинных классов, используемая для оценки качества классификации}
}

\newglossaryentry{id6}{
    name={Нейронная сеть},
    description={Вычислительная модель, вдохновлённая биологическими нейронными сетями, состоящая из слоёв взаимосвязанных узлов, обучаемых на данных}
}

\newglossaryentry{id8}{
    name={Обучение модели},
    description={Процесс настройки параметров нейронной сети путём минимизации функции потерь на обучающей выборке}
}

\newglossaryentry{id18}{
    name={Оптимизатор},
    description={Алгоритм, обновляющий параметры модели на основе вычисленных градиентов с целью минимизации функции потерь}
}

\newglossaryentry{id9}{
    name={Переобучение},
    description={Явление, при котором модель слишком точно подстраивается под обучающую выборку и теряет способность к обобщению}
}

\newglossaryentry{id3}{
    name={Транзакция},
    description={Операция передачи Bitcoin между адресами, записанная в блокчейне}
}

\newglossaryentry{id17}{
    name={Функция потерь},
    description={Функция, измеряющая расхождение между предсказаниями модели и истинными метками, минимизируемая в процессе обучения}
}

\newglossaryentry{id16}{
    name={Эпоха},
    description={Один полный проход по всему обучающему датасету в ходе обучения нейронной сети}
}

\newglossaryentry{id1}{
    name={Bitcoin},
    description={Криптовалюта и одноимённая платёжная система, использующая технологию блокчейн для децентрализованного учёта транзакций}
}

\newglossaryentry{id5}{
    name={UTXO},
    description={Unspent Transaction Output~--- модель учёта средств в Bitcoin, где каждый выход транзакции может быть потрачен только один раз}
}
